\chapter{Evaluation, Critical Appraisal and Mitigations}
\label{ch:evaluation}

This chapter asks what the preceding two chapters are worth. It separates what the
experiment demonstrates from what the report argues, reads the measured transfer in the
light of what happens to the responders as well as the forecaster, applies the definition
adopted in Chapter~\ref{ch:introduction} to the result, and then sets out at length what the
work cannot establish. It closes with mitigations and with a comparison against related
work. The conclusion of the first half is that the experiment does not settle the question it
was built to address, for reasons that are themselves informative: one limb of the
definition is satisfied, the other is undetermined because both conditions testing it turn
on a definitional problem the report inherits from the MEV literature, and the definition's
exclusion for risk compensation is set to zero by an assumption of the model. What survives is a design critique that does not depend on the verdict.

\section{What is demonstrated and what is argued}
\label{sec:demonstrated}

The distinction matters enough to be made in a table rather than left to the reader.
Table~\ref{tab:epistemic} sorts the report's substantive claims by the kind of support each
has.

\begin{table}[htbp]
\centering
\small
\begin{tabularx}{\textwidth}{p{0.30\textwidth}p{0.19\textwidth}X}
\toprule
\textbf{Claim} & \textbf{Status} & \textbf{Basis} \\
\midrule
The Sapience relayer broadcasts each request to every connected client, carrying the full pick set, the size and the predictor's address & Established & The protocol's own builder guide and implementation, read at a pinned commit \\
Disclosure moves surplus from an informed forecaster to responders & Demonstrated & Measured against a matched benchmark; the sign was expected on theory \\
The share of the forecaster's advantage lost rises with forecaster accuracy & Demonstrated & The $q_F$ sweep, every step resolved \\
Direction alone accounts for most of the transfer & Demonstrated & The $B$ sweep \\
Responders capture the whole of $L$ relative to the matched benchmark & Demonstrated & The change in responder surplus equals $L$ in all 64 cells \\
Responder surplus is the negative of forecaster surplus, and responders lose in absolute terms & True by construction & An accounting identity; the model has no uninformed order flow \\
Whether $L$ is rent or compensation & Not settled & Section~\ref{sec:howmuch}; risk is assumed away and the comparison population is absent \\
Accumulation across requests confers advantage & Argued only & Not modelled; Section~\ref{sec:limitations}(c) \\
Multi-leg requests disclose joint belief & Argued only & Not implemented \\
Semantic validity, reproducibility and attestation are distinct, and neither protocol's verification machinery guarantees the first & Argued & Chapter~\ref{ch:protocols}, from primary documentation \\
\bottomrule
\end{tabularx}
\caption{The report's claims by epistemic status. The two rows that matter most for the
overall argument, accumulation and semantic validity, are both argued rather than measured.}
\label{tab:epistemic}
\end{table}

\section{Who captures the transfer, and what the model cannot say about it}
\label{sec:responders}

Chapter~\ref{ch:results} reported $L$ as a transfer away from the forecaster. Establishing
who receives it requires care, because one natural way of asking the question gives an
answer that is true by construction and tells us nothing.

\paragraph{The uninformative version.} Responder surplus in this model is the negative of
forecaster surplus, exactly, in all sixty-four cells of the six sweeps. That is an accounting
identity, not a measurement: Chapter~\ref{ch:method} notes it, the model is zero-sum between
two parties, and the forecaster is the only trader and is always informed. It follows that
responders lose money in absolute terms in every configuration, and no parameter setting
could make them profitable. Nothing about disclosure can be inferred from that fact.

The structural reason matters and is worth stating, because it bounds the model. In Glosten
and Milgrom, a market maker survives because gains from uninformed traders offset losses to
informed ones \cite{glosten1985bidask}. This model contains no uninformed order flow at all.
Its responders therefore face a population that is informed with probability one, and their
absolute losses are a consequence of that omission rather than a finding about the
mechanism. The model is silent on what spread would make a responder solvent, and therefore
silent on what fraction of a real venue's quoted spread $L$ would represent.

\paragraph{The informative version.} Rent is defined against a counterfactual, never against
zero, and the matched no-disclosure benchmark is precisely the counterfactual the experiment
was built to supply. Measured against it, responders' surplus rises from $-0.3749$ to
$-0.1000$ at the baseline, a gain of $0.2749$. In every one of the sixty-four cells that gain
equals $L$ to the last digit. Relative to the benchmark, the responders capture the whole of
the transfer.

That is the correct reading and it points the opposite way from the absolute figures.
Disclosure does not dissipate the forecaster's advantage; it moves it, entirely, to the
responders who quote on the request. What the model cannot say is whether receiving it leaves
those parties better off than a fair price for the service they provide, because it has
removed the population against which that judgement would be made.

\section{How much of the transfer meets the definition}
\label{sec:howmuch}

Chapter~\ref{ch:introduction} defined AI-MEV as value captured because of role-specific
protocol access, beyond what is attributable to forecasting ability or to ordinary
competitive compensation for risk borne, and set out four conditions under which a transfer
might qualify. This section applies them, and takes the risk exclusion in the definition
first because it governs the rest. The result is that the experiment does not settle the question, and the reasons are
more interesting than a verdict either way would have been.

\paragraph{Compensation for risk: exactly zero within the model, and unmeasurable outside
it.} The definition excludes ordinary competitive compensation for risk borne. This model's
responders are risk-neutral, and a risk-neutral agent requires no compensation for risk. It
follows that within the model the risk component of $L$ is zero and the transfer is entirely
informational. This is the opposite of a concession, and it must not be mistaken for
evidence: the model does not measure that the transfer is not risk compensation, it assumes
a world in which no risk compensation is owed. What a risk-averse responder facing inventory
and capital constraints would require is outside the experiment entirely, and
Section~\ref{sec:limitations}(a) treats that as the deepest limitation of the work.

\paragraph{No priced option not to disclose: satisfied.} The documentation states that
unsigned requests receive no actionable bids and should be treated as price discovery only
\cite{sapience2026builderguide}. A forecaster who wants a fill must broadcast the request in
full: every pick with its direction, the size, and an address. There is no mechanism for
revealing a coarser version of the request one is actually making.

The $B$ sweep of Figure~\ref{fig:buckets} qualifies what this is worth, and the
qualification cuts both ways. A coarsening option would address the roughly five per cent of
the transfer that confidence granularity controls; the direction, which carries the other
ninety-five per cent, cannot be withheld at all by someone intending to trade. So the option the
mechanism withholds is the cheaper one. The condition is satisfied as stated, and
what it is worth to a forecaster is smaller than the phrasing suggests. That figure also
depends on the correspondence between disclosed size and confidence which
Section~\ref{sec:limitations}(e) records as unverified.

\paragraph{Asymmetric access: contested, and the contest is definitional.} The condition asks
whether some parties see order flow that others do not, or see it earlier.
Chapter~\ref{ch:protocols} established that the relayer broadcasts to every connected client
with no allowlist and no filtering \cite{sapience2026relayer,sapience2026builderguide}. It
did not establish that anyone may become a connected client, and recorded that claim as
argued rather than demonstrated; this chapter does not upgrade it. On the second limb, the
broadcast loop iterates connections in order, so some client is written to first, and
nothing in this report measures whether that ordering confers advantage.

The harder problem is definitional and the report should confront it rather than resolve it
conveniently. Suppose access is indeed open to anyone who connects. Ethereum's public
mempool is also open to anyone who connects, and participation there is gated by capital and
infrastructure rather than by permission; that is the canonical setting of the MEV
literature this project builds on \cite{daian2019flashboys,ethereummev}. A reading of
``role-specific protocol access'' under which capital-gated access to a public stream cannot
qualify would exclude most of what MEV means in practice, and would leave the framework
unable to identify the phenomenon in the setting where it is least disputed. A reading under
which it does qualify makes the condition nearly always satisfied and correspondingly weak.
This report does not resolve that tension. It notes that the definition adopted in
Chapter~\ref{ch:introduction} inherits it from the literature rather than introducing it,
and that any verdict on this condition is a verdict about the definition rather than about
Sapience.

\paragraph{Accumulation: argued, not demonstrated.} An address appears in every broadcast,
so a per-address history is cheap to build from an unprivileged position by anyone the
forecaster does not defeat by rotating. That
is documented protocol behaviour. Whether it converts into an advantage that competitive
quoting does not compensate is not measured anywhere in this report, because responders
carry no state across events and the concentration parameter was not implemented. A reader
is entitled to treat this limb as unsupported, and Section~\ref{sec:limitations}(c) lists it
as the most consequential gap.

\paragraph{Role-conditional: satisfied in form, subject to the same definitional question.}
The advantage attaches to occupying the responder position rather than to forecasting well,
which is what the condition asks. Whether occupying that position counts as protocol-specific
when it requires capital rather than permission is the question of the preceding paragraph
and has the same answer, which is that the report does not settle it.

\paragraph{The conclusion, stated plainly.} The aggregation rule of
Section~\ref{sec:definition} requires one condition on each limb. In the terms defined
there, the second limb asks whether the transfer is mechanism-mediated information rent
rather than epistemic reward, and the answer is that it is: the forecaster earns less for
the same forecasting ability solely because of how the request is broadcast. That limb
\emph{is} satisfied: the mechanism offers no priced
option not to disclose, established from the protocol's own documentation, and the
accumulation condition would reinforce it if it could be measured. The first limb, that the
access is role-specific, is undetermined, and it is undetermined for one reason rather than
four: both of its conditions turn on whether capital-gated access to a public stream counts
as protocol-specific, and Section~\ref{sec:definition} set out why that question is
inherited from the MEV literature rather than answerable here.

So the verdict fails on a single identified point, not on a vague insufficiency of evidence.
That is worth stating precisely, because it says what would settle the matter: a defensible
reading of role-specificity that survives the public-mempool objection would decide this
case immediately, given that the other limb is already satisfied. The risk exclusion in the
definition is a separate matter and is zero within the model by assumption, so it neither
blocks nor supports the verdict. \textbf{The experiment does not establish that the measured transfer is AI-MEV, and it does
not establish that it is ordinary adverse selection.} It cannot do the second because it has assumed risk away and removed the
uninformed order flow against which compensation would be judged; it cannot do the first
because the two conditions that would carry the claim are one unmeasured and one definitional.

An earlier draft of this chapter concluded that most of the transfer was ordinary adverse
selection. That conclusion rested on the absolute losses discussed in
Section~\ref{sec:responders}, which are an artefact of the missing uninformed traders, and
on reading risk-neutrality as evidence for a compensation interpretation when it implies the
reverse. It is withdrawn. Conceding a result on unsound reasoning is not more honest than
claiming one, and it would have been harder for a reader to catch.

What survives is not a verdict but a design critique, and it does not depend on the
definitional question. The mechanism compels complete, address-linked disclosure and offers
nothing cheaper. Against the matched benchmark, the whole of the forecaster's lost surplus
accrues to the responders who quote on it. And the share lost rises with the
forecaster's accuracy, from twenty-six per cent at the lowest accuracy modelled to
ninety-six per cent at the highest (Figure~\ref{fig:qf}), so the mechanism's informational
cost scales with the quality of the information supplied. Whatever that transfer is called,
a prediction market whose disclosure cost rises with forecaster skill is taxing the input it
exists to attract.

\section{Semantic validity is not reproducibility is not attestation}
\label{sec:threelayer-eval}

The second contribution is developed in Chapter~\ref{ch:protocols} and is not restated here.
Three points bear on the evaluation.

The first is that it needs no experiment, and is therefore unaffected by every limitation in
Section~\ref{sec:limitations}. The separation between the three properties, the technical
reason reproducibility and attestation are not substitutes, and the observation that neither
protocol's verification machinery guarantees semantic validity, all rest on primary
documentation read directly
\cite{gensyn2026receipts,truebit2026orchestration,uma2026oracle}.

The second is a qualification of how strongly Delphi's architectural history supports the
separation, and the weight it can bear. Chapter~\ref{ch:protocols} establishes the removal
of execution discretion and does not assert what remains, because
Section~\ref{sec:limitations}(h) records that no first-party source establishes who writes
the settlement specification under the automated-settlement deployment. What the sources
support is that a designer closed execution discretion and published nothing to indicate
that specification discretion was also closed. That is evidence, because a protocol that had
removed specification power would have reason to say so, but it is an argument from silence
and the silence has a competing explanation, since Gensyn currently document no settlement
mechanism at all. Every chapter that cites the history carries that weaker reading.

The third is that this contribution is the more durable of the two, precisely because it
does not depend on the definitional question that Section~\ref{sec:howmuch} could not
resolve. Whether a disclosure transfer counts as MEV is contested. Whether a verification
stack establishes that a specification was neutral is not: it does not, on the vendors' own
statements.

\section{Critical limitations}
\label{sec:limitations}

This section is deliberately the longest in the chapter and the least flattering. The
limitations are ordered from the most fundamental to the most contingent.

\paragraph{(a) The model cannot separate information rent from risk compensation, which is
the distinction the report's definition turns on.} Responders are risk-neutral and quote
their own posteriors without shading, so within the model the risk component is zero by
assumption and $L$ is entirely informational. The model therefore cannot measure what a risk-averse responder
would require, and every statement in Section~\ref{sec:howmuch} about compensation rests on
reasoning about the model's structure rather than on any measurement. This is the deepest
limitation and no amount of further sweeping would address it; it would require a model in
which responders price risk, hold inventory and choose whether to participate.

\paragraph{(b) There is no uninformed order flow, so responders cannot be solvent.} The
model's single trader is informed with probability one. In Glosten and Milgrom a market
maker breaks even because gains from uninformed traders offset losses to informed ones
\cite{glosten1985bidask}; with that population removed, responders lose in every
configuration and no parameter setting could change it. A second mechanism compounds this:
because the forecaster takes the best of $n$ quotes, the responder that trades is the one
whose signal was most misleading in the forecaster's favour, and the model gives it no way to
shade against that selection. This is why the absolute surplus
figures carry no information, and why Section~\ref{sec:responders} reads the transfer
against the matched benchmark instead. It also means the model has no participation
constraint and cannot say what spread would make a responder willing to quote, so its
magnitudes are upper bounds on the forecaster's position rather than estimates of any real
one.

The consequence for interpretation is set out at the point the benchmark is chosen, in
Section~\ref{sec:benchmark}, and is repeated here because it bounds every number in
Chapter~\ref{ch:results}: measured against a benchmark no responder would join, $L$ is as
defensible read as the share of a structural loss that disclosure recovers as read as
surplus transferred. The experiment does not distinguish the two.

\paragraph{(c) Accumulation, the limb that most needs support, has none.} The report's
argument that an address in every broadcast confers a durable advantage is structural and
unmeasured. Responders carry no state across events. Whether an observer who
builds a per-forecaster history can price better than competition permits, and whether the
resulting advantage would survive a forecaster who rotated addresses, are both open. This is
the single place where a reader who wished to reject the report's argument would find the
least resistance.

\paragraph{(d) One event, one forecaster, one round, and no decision to trade.} There is no
repeated play, no reputation, no competition between forecasters and no option to withhold a
request when the expected disclosure cost exceeds the expected gain. A forecaster who could
decline faces a different problem from the one modelled, and the ability to decline is
plausibly the most important strategic response available to a real participant.

\paragraph{(e) The disclosure dial is a modelling device, not a feature of either venue.}
$B$ coarsens confidence into bands. Sapience broadcasts an exact size. The correspondence
argued in Chapter~\ref{ch:method}, that size stands in for confidence because a more
confident forecaster stakes more, is plausible and unverified; no data on the relationship
between forecaster confidence and position size at any venue is used anywhere in this
report. The $B$ sweep should be read as characterising a family of mechanisms rather than as
measuring this one.

\paragraph{(f) Traded size is constant, which removes a real strategic choice.} Holding size
fixed was necessary to keep $B$ a clean disclosure dial, and it removes the forecaster's
ability to reduce disclosure by trading smaller. A forecaster who split a position across
several requests would disclose less per request at the cost of executing in pieces, and
that trade-off is not modelled.

\paragraph{(g) Multi-leg disclosure is argued and never implemented.} The strongest form of
the disclosure argument, that a combination reveals joint rather than marginal belief, rests
entirely on Chapter~\ref{ch:protocols}'s reading of the payload. No experiment supports it.

\paragraph{(h) The protocol analysis describes a moving target, and one half of the
comparison may be out of date.} Chapter~\ref{ch:protocols} could not establish whether
prompt authorship and model selection remain with the market creator under Delphi's
automated-settlement deployment, because Gensyn's current documentation describes no
settlement mechanism. The Delphi half of the resolution comparison therefore describes a
design documented in May 2026 that may no longer be deployed, while the Sapience half
describes current behaviour. That asymmetry is unavoidable given the sources and it weakens
the comparison.

\paragraph{(i) The attested-settlement architecture is not established as Delphi's.} Truebit
describe an unnamed client. The report's evolution argument is constructed not to depend on
the identification, but a reader who assumed the identification would draw a stronger
conclusion than the sources support.

\paragraph{(j) Reported numbers come from one seed, though the conclusions do not depend on
it.} Every figure in Chapter~\ref{ch:results} is drawn from a single realisation, and the
paired intervals capture Monte Carlo error within that realisation rather than across seeds.
Section~\ref{sec:res-seeds} reports a check at five seeds: every qualitative claim holds at
all five, and the headline statistics move by less than one percentage point. The residual
limitation is narrower than it was but is not zero, since five seeds bound variation rather
than eliminating it, and the check was run on the sweeps as specified rather than on
alternative specifications of the model.

\paragraph{(k) No live-venue data of any kind.} Nothing in this report measures anything
that happened on Delphi or Sapience. The ethical boundary described in
Chapter~\ref{ch:introduction} was held, and the cost of holding it is that every empirical
claim is about a simulator.

\section{Mitigations}
\label{sec:mitigations}

The mitigations below follow from the measurements where they can and are argued where they
cannot. Each is assessed against what Chapter~\ref{ch:results} showed.

\paragraph{Coarsening the disclosed size.} A venue could broadcast a size band rather than an
exact size. The $B$ sweep prices this directly and the answer is discouraging: moving from
sixty-four bands to one recovers about $5\%$ of the transfer, because direction alone
accounts for the rest. Coarsening size is cheap and nearly useless.

\paragraph{Withholding direction.} Since direction carries $95\%$ of the effect, a mechanism
that solicited quotes without revealing which side the requester wanted would address the
part that matters. Two-sided quoting is the standard instrument, and it must not be confused
with this experiment's no-disclosure arm: in that arm responders never learn the direction
at all, whereas under two-sided quoting the responder posts both sides and the informed
forecaster then chooses which to hit, which is the Glosten-Milgrom setting in which the
market maker is picked off \cite{glosten1985bidask}. Two-sided quoting does not eliminate
adverse selection; it prices it into the spread. This experiment does not model
side-selection against a posted quote and so cannot say how much of the $95\%$ such a
mechanism would recover. The measurement identifies direction as the quantity worth
protecting; which instrument protects it best it cannot settle.

\paragraph{Sealed bids and delayed reveal.} Committing bids before the request is revealed,
or revealing the request only to the winning responder, would approximate the no-disclosure
arm. This is argued rather than measured: the experiment models what responders know, not
the mechanism by which they come to know it, so it cannot distinguish a sealed-bid scheme
from any other means of achieving the same information state.

\paragraph{Noised or aggregate disclosure.} Publishing a differentially private hint rather
than the request itself is the approach taken in order-flow auctions
\cite{passeratpalmbach2025dphints}, and the closest analytic treatment derives a closed-form
spread and a per-trade welfare transfer for a market maker observing a noisy direction signal
in the Glosten-Milgrom setting \cite{nakamura2026privacysubsidy}, with a continuous-time
counterpart \cite{nakamura2026kyle}. Both are unrefereed preprints and both model a single
committed market maker rather than competing responders, so they bound a different problem
from the one measured here. What this report adds is a caution rather than a
result: because $I$ is invariant by construction across the responder-count and dependence
sweeps while $L$ varies by up to a quarter, a bound on how much a request reveals does not
determine the economic transfer. That is a logical consequence of the metric's definition
rather than a measurement, as Section~\ref{sec:res-info} states. It bears on the mitigation
nonetheless: a disclosure bound constrains $L$ without pinning it down, and the competitive
environment has to be modelled alongside it.

\paragraph{Evaluation inside a trusted execution environment with conditional recall.} A
responder could evaluate a request inside an enclave and be required to discard it if no
trade results, so that a request previews information without the responder retaining it.
The idea was raised as credible forgetting at the MEV-SBC agentic MEV breakout
\cite{mevsbc2026breakout}. It targets the accumulation limb, which is the limb this report
could not measure. One objection must be stated with it: a responder has to emit a quote,
and a quote conditioned on the request already encodes information about it, to the
forecaster and to anyone observing quotes. An enclave can stop a responder retaining a
request; it cannot unsay what the responder's own price reveals. Conditional recall
therefore bounds accumulation across requests rather than disclosure within one. That is
still the right target for the gap this report leaves open, but it is a narrower claim than
the idea first suggests. It is argued here and not evaluated.

\paragraph{On the resolution side, there is no mitigation of the same kind.} Nothing in the
verification stack addresses specification neutrality, so the available responses are
procedural rather than cryptographic: separating the party who writes the settlement
criterion from the party who holds a position, publishing criteria before trading opens so
that they can be contested while contesting is still cheap, and treating a settlement
specification as a governance artefact rather than a configuration field. None of these is
verifiable in the sense the rest of the stack is verifiable, which is the point.

\section{Comparison with related work}
\label{sec:related}

Table~\ref{tab:related} positions the work.

\begin{table}[htbp]
\centering
\small
\begin{tabularx}{\textwidth}{p{0.26\textwidth}p{0.28\textwidth}X}
\toprule
\textbf{Approach} & \textbf{What it establishes} & \textbf{What it does not} \\
\midrule
Classical microstructure \cite{glosten1985bidask,kyle1985continuous} & That informed order flow forces a spread, analytically and in general & Nothing specific to a broadcast RFQ, a public order-flow stream, or belief as the traded object \\
Order-flow auction privacy \cite{passeratpalmbach2025dphints} & That disclosure can be quantified and priced, with a mechanism for doing so & It treats transaction order flow, not forecasts, and does not address resolution \\
Noisy-observation microstructure \cite{nakamura2026privacysubsidy,nakamura2026kyle} & Closed-form spread and welfare under a noisy direction signal & A single committed market maker, not $n$ competing responders with dependent signals; unrefereed \\
Prediction-market arbitrage \cite{saguillo2025forest,gebele2026semantic} & That event semantics and fragmentation create exploitable inconsistency, within a platform and across platforms respectively & It concerns pricing rather than disclosure at execution or the settlement procedure \\
This report & A matched-benchmark measurement of disclosure under competing responders, and a separation of three verification properties & It measures no live venue, cannot separate risk from information, and does not model accumulation \\
\bottomrule
\end{tabularx}
\caption{Position of this work. The final row's third column is the honest summary of the
preceding sections: what this report adds is a controlled comparison and an analytical
separation, not a measurement of extraction.}
\label{tab:related}
\end{table}

\section{What this chapter established}

The experiment demonstrates that disclosure moves surplus away from an informed forecaster,
that measured against the matched benchmark the responders capture the whole of it, that
direction carries about ninety-five per cent of the effect, and that the share lost rises
with the forecaster's accuracy to ninety-six per cent at the highest accuracy modelled.

It does not establish whether that transfer is AI-MEV. Under the aggregation rule of
Section~\ref{sec:definition} the second limb of the definition is satisfied and the first is
undetermined, because both of the conditions testing it turn on a question this report
inherits from the MEV literature and does not resolve. The definition's risk exclusion is
set to zero by an assumption of the model rather than tested, so it neither blocks the
verdict nor supports it. An
earlier draft of this chapter concluded that most of the transfer was ordinary adverse
selection. That conclusion rested on an accounting identity and on a backwards reading of
risk-neutrality, and Section~\ref{sec:howmuch} withdraws it.

What survives without depending on the verdict is a design critique. The mechanism compels
complete, address-linked disclosure and offers nothing cheaper; the whole of the forecaster's
lost surplus accrues to the responders who quote on it; and the cost rises with the quality
of the information disclosed.

The argued contribution stands independently of all of it. Semantic validity,
reproducibility and execution attestation are three properties, neither protocol's verification
machinery guarantees the first, and the vendors say so themselves.
Chapter~\ref{ch:conclusions} takes up what follows.
